\par\noindent\textbf{(1)} 将磁场项平均分给相邻键，转移矩阵 \(T_{ss'}=\exp[\beta Jss'+\beta h(s+s')/2]\)，即 \(\begin{pmatrix}e^{\beta(J+h)}&e^{-\beta J}\\e^{-\beta J}&e^{\beta(J-h)}\end{pmatrix}\)。特征值为 \(\lambda_\pm=e^{\beta J}\cosh(\beta h)\pm\sqrt{e^{2\beta J}\sinh^2(\beta h)+e^{-2\beta J}}\)。
\par\noindent\textbf{(2)} 配分函数 \(Z_N=\operatorname{tr}(T^N)=\lambda_+^N+\lambda_-^N\)。热力学极限 \(f=-k_BT\ln\lambda_+\)，每自旋平均磁化 \(m=-\partial f/\partial h=\sinh(\beta h)/\sqrt{\sinh^2(\beta h)+e^{-4\beta J}}\)。在 \(h=0\) 该式对任意有限 \(T\) 给 \(m=0\)。
\par\noindent\textbf{(3)} 零场转移矩阵的本征值为 \(2\cosh(\beta J)\) 与 \(2\sinh(\beta J)\)，故 \(\langle s_0s_r\rangle=(\lambda_-/\lambda_+)^r=[\tanh(\beta J)]^r\)。写成 \(e^{-r/\xi}\)，关联长度 \(\xi^{-1}=-\ln[\tanh(\beta J)]\)，随低温指数增长。
\par\noindent\textbf{(4)} 任意 \(T>0\) 时 \(\lambda_+>|\lambda_-|\)，自由能解析、零场磁化恒为零，因此一维有限邻接铁磁Ising模型没有有限温自发铁磁相变。仅当 \(T\to0\)，\(\tanh(\beta J)\to1\)，关联长度发散；引入长程作用可能改变结论。
